% [ck: Feb  7 89]
% les macros speciales pour le livre

%%%%%%%% traitement des exercices a la Knuth:
%%%%% 1. le fichier de reponses et le compteur
\newwrite\ans
\immediate\openout\ans=reponses
\newcount\exno\global\exno=0
% 2. la macro de definition de l'enonce
\chardef\other=12
\newcommand{\exercice}[1]
   {\global\advance\exno by 1 
    {\small\noindent {\bf Exercice~\the\exno~---} #1}\ifvmode\else\newline\fi
   }
%%%%% 3.  la macros de definition de la reponse
% fait avec Eric [ck: Aug 20 93]
\outer\def\answer{\par\medbreak
  \immediate\write\ans{}
  \immediate\write\ans{\the\exno:}
  \copytoblankline}
\def\copytoblankline{\begingroup\setupcopy\copyans}
\def\setupcopy{\def\do##1{\catcode`##1=\other}\dospecials
  \catcode`\|=\other \obeylines}
{\obeylines \gdef\copyans#1
  {\def\next{#1}%
  \ifx\next\empty\let\next=\endgroup %
  \else\immediate\write\ans{\next} \let\next=\copyans\fi\next}}

%%%%% 4. Utiliser les reponses en chantant dans l'annexe A
%    \immediate\closeout\ans \input{reponses}
%%%%% 5. Utilisation typique
%       \exercice{enonce'}
%       \answer{reponse eventuelle (non obligatoire)}
%              <- ligne blanche OBLIGATOIRE apres la reponse
%                 Attention aux lignes blanches dans la reponse!
%%%%% 6. Ancienne macro
\newcommand{\exer}[2]{\global\advance\exno by 1 
    {\small\noindent {\bf Exercice~\the\exno~---} #1}\ifvmode\else\newline\fi
     }          %%%%%%% \noindent{\bf Answer}: #2}
%%%%%%%%%%%%%%%%%%%%%%%%%%%%%%%%%%%%%%%%%%%%%%%%%%%%%%%%%%%%%%%%
\newcommand{\coq}{\textsf{Coq}}
\newcommand{\elan}{\textsf{ELAN}}
\let\ELAN=\elan
\newcommand{\REF}{\textsf{REF}}

% Lecture supplementaires: Further Readings
\newcommand{\FR}[1]{\noindent{\it Further Readings: #1}}

% attention passage dangereux
\newcommand{\zz}{$\bigodot$}

\newcommand{\Helene}{H{\'e}l{\`e}ne}

% Impression d'un passage en deux colonnes
% attention a la separation des colonnes et a leur largeur

\newcommand{\deuxcol}[2]{
        \noindent
        \begin{minipage}[t]{5.5cm}
        #1
        \end{minipage} \ ~~~ \
        \begin{minipage}[t]{5.5cm}
        #2
        \end{minipage}
        }

 % pour traiter les theories equationnelles introduites dans le livre
 % et faire un index specifique des theo equa
\newcommand{\thname}[1]{\index{#1|sfac}{\sf #1}}
\newcommand{\trsname}[1]{\overrightarrow{\sf #1}} 

% pour les fleches longues une abbrev:
\newcommand{\vect}[1]{\overrightarrow{#1}}

\newcommand{\bold}[1]{{\bf #1}}
\newcommand{\ital}[1]{{\it #1}}
\newcommand{\sfac}[1]{{\sf #1}}

\newcommand{\ul}[1]{\underline{#1}}
\newcommand{\ol}[1]{\overline{#1}}

\newcommand{\subsubsubsection}[1]{{\sc #1 \break}}
\newcommand{\DEF}{\begin{definition}\begin{rm}}
\newcommand{\FDEF}{\end{rm}\end{definition}}
\newcommand{\PROP}{\begin{proposition}\begin{rm}}
\newcommand{\FPROP}{\end{rm}\end{proposition}}
\newcommand{\COR}{\begin{corollary}\begin{rm}}
\newcommand{\FCOR}{\end{rm}\end{corollary}}
\newcommand{\EX}{\begin{example}\begin{rm}}
\newcommand{\FEX}{\end{rm}\end{example}}
\renewcommand{\TH}{\begin{theorem}}
\newcommand{\FTH}{\end{theorem}}
\newcommand{\LEM}{\begin{lemma}\begin{rm}}
\newcommand{\FLEM}{\end{rm}\end{lemma}}

\newcommand{\Def}[1]{{\em #1}\index{#1}} % |bold}}
\newcommand{\Aut}[1]{{#1}\index{#1|ital}}
        % typiquement pour indexer les auteurs cite's \Aut{Robinson}


\newcommand{\exligne}[1]{\begin{quote}{\tt #1} \end{quote}}
\newcommand{\Not}[1]{#1}
        % typiquement pre'vu pour pouvoir creer un index des notations
%\newcommand{\comment}[1]{**\marginpar{**{\sf\small #1}}}
\newcommand{\comment}[1]{\footnote{{\bf CHK:} #1}}
        % typiquement \comment{commentaire qui n'apparaitra pas dans
        % la version finale}
\newcommand{\apply}[2]{#1(#2)}


\newcommand{\proc}{ {\bf proc} }
\newcommand{\returns}{ {\bf returns} }
\newcommand{\signals}{ {\bf signals} }
\newcommand{\return}{ {\bf return} }
\newcommand{\signal}{ {\bf signal} }
\newcommand{\resignal}{ {\bf resignal} }
\newcommand{\si}{ {\bf if} }
\newcommand{\alors}{ {\bf then} }
\newcommand{\sinon}{ {\bf else} }
\newcommand{\fin}{ {\bf end} }
\newcommand{\fsi}{ {\bf endif} }
\newcommand{\pour}{ {\bf for} }
\newcommand{\faire}{ {\bf do} }
\newcommand{\dans}{ {\bf in} }
\newcommand{\finpour}{ {\bf enddo} }
\newcommand{\jusqua}{ {\bf to} }
\newcommand{\cas}{ {\bf case} }
\newcommand{\fincas}{ {\bf endcase} }

 % pour les exemples

\newcommand{\succe}{s}
\newcommand{\prede}{p}

\newcommand{\equi}{\rightleftharpoons}
\newcommand{\equiA}{\rightleftharpoons_{\cal A}}

\newcommand{\repla}{\hookleftarrow} 
\newcommand{\simplif}{\hookrightarrow} 
    % pour les preuves inductives par reecriture
\newcommand{\la}{\leftarrow}
\newcommand{\rra}{\rightarrow\!\!\!\!\rightarrow}
\newcommand{\lla}{\leftarrow\!\!\!\!\leftarrow}
\newcommand{\rrra}{\rightarrow\!\!\!\!\rightarrow\!\!\!\!\rightarrow}
\newcommand{\da}[1]{\downarrow_{#1}}
%\newcommand{\longla}{\longleftarrow}
\newcommand{\La}{\Leftarrow}
\newcommand{\Lra}{\Leftrightarrow}
\newcommand{\Ra}{\Rightarrow}
%-> notation \newcommand{\ra}{\rightarrow}
%\newcommand{\Longla}{\Longleftarrow}
\newcommand{\Longra}{\Longrightarrow}
%\newcommand{\Longlra}{\Longleftrightarrow}
%\newcommand{\Longra}[1]{\stackrel{#1}{\Longrightarrow}}
\newcommand{\Longla}[1]{\stackrel{#1}{\Longleftarrow}}
\newcommand{\Longlra}[1]{\stackrel{#1}{\Longleftrightarrow}}
%\newcommand{\longra}{\longrightarrow}
\newcommand{\longlra}{\longleftrightarrow}
\newcommand{\rhu}{\rightharpoonup}
\newcommand{\rlh}{\rightleftharpoons}
\newcommand{\lrh}{\rightleftharpoons}
\newcommand{\unequiv}{\vdash\!\dashv}
\newcommand{\equiva}[2]{\vdash\!#1\!\dashv_{[#2]}}
\newcommand{\longra}[1]{\stackrel{#1}{\longrightarrow}}
\newcommand{\longla}[1]{\stackrel{#1}{\longleftarrow}}
\newcommand{\lra}[1]{\stackrel{#1}{\longleftrightarrow}}
\newcommand{\reec}[3]{\stackrel{#1}{\longrightarrow}^{#2}_{#3}}
        % typiquement: \reec{*}{RE}{[3]} pour -*->RE[3] [ck: May 29 88]

% pour la reecriture conditionelle
%\newcommand{\condeq}[3]{{#1} \; \Rightarrow \;{#2} = {#3}}
%\newcommand{\cond}[3]{{#1} \; \Rightarrow \;{#2} \rightarrow {#3}}
%\newcommand{\horn}[2]{{#1} \Rightarrow {#2}}
%\newcommand{\Grd}[1]{Ground{#1}} % a modifier peut-etre
%\newcommand{\ord}{\succ}
%\newcommand{\ordC}{\succ_C}
%\newcommand{\ordE}{\succ_E}
%\newcommand{\ordeq}{\succeq}
\newcommand{\ordeqC}{\succeq_C}
\newcommand{\ordeqE}{\succeq_E}

\newcommand{\meta}[1]{\dot{#1}}


\newcommand{\hligne}{\rule{\textwidth}{0.5pt}}



%%%% pour avoir une numerotation x.y ou x est le numero de section ou
%%%% de chapitre suivant le type de document. [ck: Jan 31 94] + Eric

\edef\maxunit{\ifx\chapter\undefined section\else chapter\fi}

\ifx\theorem\undefined
\newtheorem{theorem}{Theorem}[\maxunit]
\fi
\ifx\proposition\undefined
\newtheorem{proposition}{Proposition}[\maxunit]
\fi
\ifx\lemma\undefined
\newtheorem{lemma}{Lemma}[\maxunit]
\fi
\ifx\example\undefined
\newtheorem{example}{Example}[\maxunit]
\fi
\ifx\definition\undefined
\newtheorem{definition}{Definition}[\maxunit]
\fi
\ifx\corollary\undefined
\newtheorem{corollary}{Corollary}[\maxunit]
\fi

\newcommand{\NOT}{\noindent {\bf Notation: }}
\newcommand{\FNOT}{}
\newcommand{\EXER}{\noindent {\bf Exercise: }}
\newcommand{\FEXER}{}
\newcommand{\REM}{\noindent {\bf Remark: }}
\newcommand{\FREM}{}

\ifx\proof\undefined
\newcommand{\proof}[1]{\begin{description}
            \item[Proof:] #1 \(\Box\) \end{description}}
\fi

\newcommand{\Note}[1]{{\bf Note:} #1}

%%%%%%%%%%%%%%%%% Les regles d'inferences

%%%% le symbole de transformation
\newcommand{\fleche}{\mathop{\mapsto\!\!\!\!\!\mapsto}}
\newcommand{\LaFleche}{\mathop{\mapsto\!\!\!\!\!\mapsto}}


%%%% Pour imprimer des ensembles de regles
% la definition de l'environnement

\newenvironment{ruleset}%
{\[\begin{array}{llcl}}{\end{array}\]}%

% la definition des re`gles a utiliser dans un environnement {ruleset}

\newcommand{\regle}[3]{
           {\mbox{\boldmath\bf #1}}     & {#2}  & \fleche       & {#3}\\}

\newcommand{\cregle}[4]{
           {\mbox{\boldmath\bf #1}}     & {#2}  & \fleche       & {#3}          \\
                        &       &                       & {\sf if}~ {#4}\\}
\newcommand{\wregle}[4]{
           {\mbox{\boldmath\bf #1}}     & {#2}  & \fleche       & {#3}          \\
                        &       &                       & {\sf where}~ {#4}\\}

\newcommand{\cwregle}[5]{
           {\mbox{\boldmath\bf #1}}     & {#2}  & \fleche       & {#3}          \\
                        &       &                       & {\sf if}~ {#4}\\
                        &       &                       & {\sf where}~ {#5}\\}

% juste une regle: [ck: Mar 26 93]

\newcommand{\uneregle}[3]{\begin{ruleset}\regle{#1}{#2}{#3}\end{ruleset}}
\newcommand{\unecregle}[4]{\begin{ruleset}\cregle{#1}{#2}{#3}{#4}\end{ruleset}}
\newcommand{\unegregle}[3]{\begin{gruleset}\gregle{#1}{#2}{#3}\end{gruleset}}


%%%% Pour imprimer des ensembles de regles moyennes
% la definition de l'environnement

\newenvironment{mruleset}
               {\[ \begin{array}{lll}}{\end{array} \] }

\newcommand{\mregle}[3]{
           {\bf #1}                     & {#2}  \\
           \multicolumn{1}{r}{\fleche}  & {#3}  \\}

\newcommand{\mcregle}[4]{
           {\bf #1}     & {#2}  \\
           \multicolumn{1}{r}{\fleche}  & {#3}  & {\sf if}~ {#4}\\}
\newcommand{\mwregle}[4]{
           {\bf #1}     & {#2}  \\
           \multicolumn{1}{r}{\fleche}  & {#3}  & {\sf where}~ {#4}\\}

% pour causer, en mode texte, du nom des regles
\newcommand{\rname}[1]{\mbox{\boldmath\bf #1}}

%%%% Pour imprimer des ensembles de grosses regles
% la definition de l'environnement

\newenvironment{gruleset}
               {\[ \begin{array}{ll}}{\end{array} \] }

% la definition des re`gles a utiliser dans un environnement {gruleset}

\newcommand{\gregle}[3]{
           {\mbox{\boldmath\bf #1}}     & {#2}                    \\
                        & \fleche         \\
                        & {#3}                    \\}


\newcommand{\gcregle}[4]{
           {\mbox{\boldmath\bf #1}}     & {#2}                    \\
                        & \fleche         \\
                        & {#3}                    \\
                        & {\sf if}~ {#4}          \\}

\newcommand{\gwregle}[4]{
           {\mbox{\boldmath\bf #1}}     & {#2}                    \\
                        & \fleche         \\
                        & {#3}                    \\
                        & {\sf where}~ {#4}          \\}
\newcommand{\gcwregle}[5]{
           {\mbox{\boldmath\bf #1}}     & {#2}                    \\
                        & \fleche         \\
                        & {#3}                    \\
                        & {\sf if}~ {#4}          \\
                        & {\sf where}~ {#5}          \\}

%%%% Pour imprimer des ensembles de Tres Grosses regles
% la definition de l'environnement

\newenvironment{tgruleset}
               {\[ \begin{array}{l}}{\end{array} \] }

% la definition des re`gles a utiliser dans un environnement {tgruleset}

\newcommand{\tgregle}[3]{
           {\bf #1}     \\
           {#2}         \\
           \fleche        \\
           {#3}     \\   }
 
\newcommand{\tgcregle}[4]{
           {\bf #1}     \\ 
           {#2}         \\
           \fleche        \\
           {#3}         \\
           {\rm if}~ {#4}\\  }
